\chapter{La Geometria: l'Apparato
Macro}\label{la-geometria-lapparato-macro}

\begin{quote}
\textbf{Argumento central:} ogni attore egemonico occupa una posizione
in uno spazio tridimensionale di stato (coesione, capacità,
legitimidade); la sua interazione strategica con altri attori segue un
payoff la cui forma è identica a quella del singolo acoplamento
discorsivo della Topografia; la stabilità del sistema egemonico è una
soglia \(T_s\) che misura il rapporto tra forze di coerenza e di
decoerência.
\end{quote}



\section{Perché un Capítulo Sulla
Geometria}\label{perchuxe9-un-capítulo-sulla-geometria}

Il Cap.~1 ha presentato la Topografia come microfundação comunicativa.
Questo capítulo presenta l'apparato simmetrico al livello
macro-geopolitico: la \emph{Geometria dell'Egemonia} (Marcelino, 2026),
che applica formalismo geometrico, teoria dei giochi e estensione
stocastica alla dinamica tra attori egemonici. Il Cap.~3 dimostrerà che
le due strutture sono \emph{unificabili} através de l'operatore di
Kuramoto. Per seguire quella demonstração, o leitor ha bisogno
dell'apparato della Geometria --- questo capítulo lo fornisce in forma
compatta.



\section{Il Problema Egemonico
Classico}\label{il-problema-egemonico-classico}

La tradizione critica ha trattato l'egemonia in registro qualitativo:
Gramsci (egemonia come direzione morale e culturale), Cox (egemonia come
articolazione di forze materiali, idee, istituzioni), Wallerstein (cicli
sistemici). Manca un \textbf{formalismo unificato} che permetta
predizioni datate e falsificabili. La Geometria propone uno spazio
metrico delle posizioni egemoniche e una dinamica formale delle loro
interazioni.



\section{\texorpdfstring{Lo Spazio degli Stati Egemonici
\(\mathcal{H} = [0, 10]^3\)}{Lo Spazio degli Stati Egemonici \textbackslash mathcal\{H\} = {[}0, 10{]}\^{}3}}\label{lo-spazio-degli-stati-egemonici-mathcalh-0-103}

Ogni attore egemonico \(i\) è caratterizzato da uno stato
\(\mathbf{s}_i = (s_{i,1}, s_{i,2}, s_{i,3}) \in \mathcal{H}\):

\begin{itemize}
\tightlist
\item
  \(s_{i,1} = s_i\) = \textbf{coesão interna} (integrazione politica,
  omogeneità decisionale, controllo del territorio);
\item
  \(s_{i,2} = s_{ii}\) = \textbf{capacidade material} (PIL, popolazione,
  tecnologia, energia, forze armate);
\item
  \(s_{i,3} = s_{iii}\) = \textbf{legitimidade ético-discursiva} (soft
  power, legitimidade internazionale, qualità del discorso pubblico).
\end{itemize}

\Hypoth{} \textbf{Disambiguazione notazionale.} Lo spazio
\(\mathcal{H}\) degli stati egemonici è denotato con la lettera \emph{H}
(per \emph{Hegemonic}) per evitare conflitto con \(\mathcal{S}\) del
Cap.~5, che indica la \emph{superficie critica}
\(T_s^{(2),\text{sys}} = 1\) in \(\mathbb{V}^{(4)}\). I due oggetti
hanno natura matematica e ruolo diversi: \(\mathcal{H}\) è uno spazio di
stato statico tridimensionale; \(\mathcal{S}\) è un'ipersurface
(sottovarietà di codimensione 1) nello spazio estrutural
quadridimensionale.

\Hypoth{} \textbf{Mappa con la Topografia.} Gli assi
\((s_i, s_{ii}, s_{iii})\) sono gli omologhi di scala degli assi
\((x, y, z)\) della Topografia. La calibrazione empirica di
\(\Sigma^{\text{geo}}\) a partir de \(\Sigma^{\text{comm}}\) è discussa
nell'Apêndice B della DEQ$^2$.

\subsection{La Qualità Strategica}\label{la-qualituxe0-strategica}

\[
Q_i = \frac{s_{i,1} + s_{i,2} + s_{i,3}}{\sqrt{3}}
\]

è la qualità complessiva dell'attore \(i\) --- formalmente identica al
\(Q_0\) della Topografia, applicata agli stati egemonici em vez disso che alle
componenti del discorso.



\section{Il Payoff Egemonico e il Folk
Theorem}\label{il-payoff-egemonico-e-il-folk-theorem}

\subsection{Payoff di Periodo}\label{payoff-di-periodo}

L'attore \(i\) ottiene, per ogni periodo, un payoff:

\[
\Pi_i(t+1) = H_i \cdot Q_i \cdot (1 - \delta_i \cdot s_{i,3})
\]

onde \(H_i\) è l'\textbf{impronta sistemica} (la ``massa egemonica'' ---
analogo macro di \(M_e\)) e \(\delta_i\) il fator de desconto. Il
termine \((1 - \delta_i s_{i,3})\) è il \textbf{fator de erosão di
legitimidade}: la legitimidade conquistata si consuma nell'esercizio del
potere.

\Proved{} \textbf{Identità formale con Topografia (verificaçãota nel Cap.~3
\S{}3.1.3).} Il payoff egemonico è strutturalmente la stessa equação
della função preditiva della Topografia
\(\Pi(t+1) = M_e Q_0 (1 - \delta_R z_R)\), con corrispondenza:

{\def\LTcaptype{none} % do not increment counter
{\small\begin{longtable}[]{@{}
  >{\raggedright\arraybackslash}p{(\linewidth - 2\tabcolsep) * \real{0.5000}}
  >{\raggedright\arraybackslash}p{(\linewidth - 2\tabcolsep) * \real{0.5000}}@{}}
\toprule\noalign{}
\begin{minipage}[b]{\linewidth}\raggedright
Topografia
\end{minipage} & \begin{minipage}[b]{\linewidth}\raggedright
Geometria
\end{minipage} \\
\midrule\noalign{}
\endhead
\bottomrule\noalign{}
\endlastfoot
\(M_e\) (massa comunicativa) & \(H_i\) (impronta sistemica) \\
\(Q_0\) (qualidade ética) & \(Q_i\) (qualità strategica) \\
\(\delta_R\) (sconto del ricevente) & \(\delta_i\) (sconto
dell'attore) \\
\(z_R\) (complessità del ricevente) & \(s_{i,3}\) (legitimidade erosa
nell'uso) \\
\end{longtable}}
}

\subsection{Il Teorema Folk della
Geometria}\label{il-folk-theorem-della-geometria}

\Proved{} Per il gioco egemonico ripetuto, la cooperazione tra egemoni
(anziché conflitto totale) è sostenibile come equilíbrio de Nash
perfetto nei sottogiochi se e solo se:

\[
\delta_i \geq \delta^* = 0{,}50
\]

\Hypoth{} \textbf{Identità con il Teorema Folk Discursivo della
Topografia.} La stessa \(\delta^* = 0{,}50\) del Folk Discorsivo per
l'Emittente. Questa coincidenza numerica è una firma matematica di
invariância de escala --- la stessa struttura governa cooperazione
comunicativa micro e cooperazione egemonica macro. La DEQ$^2$ la dimostra
come teorema estrutural (Cap.~3).

\subsection{Implicazione Geopolitica del Folk
Theorem}\label{implicazione-geopolitica-del-folk-theorem}

\Hypoth{} Sistemi con \(\delta < 0{,}50\) --- basso horizonte temporal,
alta preferenza per il presente --- non possono cooperare stabilmente.
Il populismo del \emph{short term} (cicli elettorali brevi, retorica
trimestrale dei mercati, \emph{engagement} algoritmico) è
strutturalmente incompatibile con la cooperazione egemonica. Predizione:
i sistemi politici che riducono \(\delta\) collettivo sotto \(0{,}50\)
entrano in regime di non-cooperazione estrutural, indipendentemente
dalla volontà degli attori.



\section{\texorpdfstring{La Soglia di Sgretolamento \(T_s\)
(DEQ$^1$)}{La Soglia di Sgretolamento T\_s (DEQ$^1$)}}\label{la-soglia-di-sgretolamento-t_s-dequxb9}

Il \emph{Memoriale Tecnico DEQ} (2026-04-23) deriva, dalle equações del
payoff egemonico in regime stocastico, una soglia composita:

\[
T_s^{(1)} = \frac{s_{iii} \cdot z \cdot \delta}{\sigma \cdot \varepsilon_{\text{dis}}}
\]

onde:

\begin{itemize}
\tightlist
\item
  \(s_{iii}\) = legitimidade ético-discursiva del sistema;
\item
  \(z\) = proiezione sull'asse dialettico (eredità diretta della
  Topografia);
\item
  \(\delta\) = horizonte temporal (Teorema Folk);
\item
  \(\sigma\) = volatilidade estocástica ambientale (estensione EDSD, \S{}2.5);
\item
  \(\varepsilon_{\text{dis}}\) = coeficiente de desengajamento moral
  (Bandura 1999, Milgram 1974).
\end{itemize}

\Proved{} \textbf{Lettura.} Numeratore: forze di coerenza. Denominatore:
forze di decoerência. La limiar crítico è \(T_s^{(1)} = 1\). Sopra,
sistema stabile; sotto, transição de fase imminente.

\Conj{} \textbf{Limite della DEQ$^1$.} L'equação \emph{non distingue} tra
sistemi resilienti e antifragili, né tra attrito esogeno (volatilità) ed
endogeno (involução). Due sistemi con stesso \(T_s^{(1)}\) possono
avere traiettorie radicalmente diverse. Questa lacuna motiva
l'estensione DEQ$^2$ con \(A\) e \(\Omega\) (Cap.~3).



\section{L'Estensione Stocastica
EDSD}\label{lestensione-stocastica-edsd}

La \emph{Geometria dell'Egemonia} (Cap.~12-13) introduce
un'\textbf{Equazione Dinamica Stocastica del Decadimento} che modella
l'evoluzione temporale di ciascuno stato \(s_{i,k}\) come processo
stocastico:

\[
ds_{i,k} = \mu_k(\mathbf{s}_i, t)\, dt + \sigma_k(\mathbf{s}_i, t)\, dW_k(t)
\]

onde \(W_k(t)\) è un processo di Wiener e \(\sigma_k\) è la volatilità
sull'asse \(k\). La \(\sigma\) che entra in \(T_s^{(1)}\) è una
\textbf{norma aggregata} del vettore \((\sigma_1, \sigma_2, \sigma_3)\),
in genere \(\sigma = \|\boldsymbol{\sigma}\|_2\).

\Hypoth{} \textbf{Conseguenza per la DEQ$^2$.} L'antifragilidade \(A\) del
Cap.~3 è la \emph{convessità della risposta} a questa \(\sigma\) ---
collega EDSD alla formalizzazione di Taleb. Sistemi con \(A > 0\)
trasformano la varianza stocastica in opzionalità.



\section{Cosa la Geometria Porta alla
DEQ$^2$}\label{cosa-la-geometria-porta-alla-dequxb2}

{\def\LTcaptype{none} % do not increment counter
{\small\begin{longtable}[]{@{}
  >{\raggedright\arraybackslash}p{(\linewidth - 4\tabcolsep) * \real{0.3333}}
  >{\raggedright\arraybackslash}p{(\linewidth - 4\tabcolsep) * \real{0.3333}}
  >{\raggedright\arraybackslash}p{(\linewidth - 4\tabcolsep) * \real{0.3333}}@{}}
\toprule\noalign{}
\begin{minipage}[b]{\linewidth}\raggedright
Apparato Geometrico
\end{minipage} & \begin{minipage}[b]{\linewidth}\raggedright
Ruolo nella DEQ$^2$
\end{minipage} & \begin{minipage}[b]{\linewidth}\raggedright
Referência
\end{minipage} \\
\midrule\noalign{}
\endhead
\bottomrule\noalign{}
\endlastfoot
Spazio \(\mathcal{H} = [0,10]^3\) degli stati egemonici & Spazio degli
stati macro & Cap.~3 \\
\(Q_i\) qualità strategica & Forma identica a \(Q_0\) Topografia
(invariância de escala) & Cap.~3 \S{}3.1.3 \\
\(\Pi_i(t+1) = H_i Q_i (1 - \delta_i s_{i,3})\) & Legge macro-payoff
identica a $\Pi$ Topografia & Cap.~3-4 \\
Teorema Folk \(\delta^* = 0{,}50\) & Stessa soglia di Folk Discorsivo
--- invariância de escala & Cap.~3 \S{}3.1.3 \\
\(T_s^{(1)} = s_{iii} z \delta / (\sigma \varepsilon_{\text{dis}})\) &
Caso limite della DEQ$^2$ per \(A = \Omega = 0\) & Cap.~3 \\
EDSD: \(ds = \mu\, dt + \sigma\, dW\) & Definizione operativa di
\(\sigma\) in \(T_s\) & Cap.~4 \\
Coefficiente di disimpegno \(\varepsilon_{\text{dis}}\)
(Bandura/Milgram) & Variabile macro che assorbe la segnalazione
bayesiana di Topografia & Cap.~3 \\
\end{longtable}}
}

\Hypoth{} \textbf{Sintesi epistemica.} Se la Topografia è la
termodinamica statistica del sistema discorsivo, la Geometria è la
\textbf{termodinamica fenomenologica} del sistema egemonico --- descrive
le leggi di trasformazione macroscopiche senza esplicitarne la
microfundação. La DEQ$^2$ (Cap.~3-4) compie il passo che unifica i due
livelli, dimostrando che le leggi macro emergono per trasposizione di
Kuramoto dalle leggi micro.



\section{Notas Epistêmicas de Encerramento
Capítulo}\label{note-epistemiche-di-chiusura-capítulo}

{\def\LTcaptype{none} % do not increment counter
{\small\begin{longtable}[]{@{}
  >{\raggedright\arraybackslash}p{(\linewidth - 6\tabcolsep) * \real{0.2500}}
  >{\raggedright\arraybackslash}p{(\linewidth - 6\tabcolsep) * \real{0.2500}}
  >{\raggedright\arraybackslash}p{(\linewidth - 6\tabcolsep) * \real{0.2500}}
  >{\raggedright\arraybackslash}p{(\linewidth - 6\tabcolsep) * \real{0.2500}}@{}}
\toprule\noalign{}
\begin{minipage}[b]{\linewidth}\raggedright
\#
\end{minipage} & \begin{minipage}[b]{\linewidth}\raggedright
Afirmação
\end{minipage} & \begin{minipage}[b]{\linewidth}\raggedright
Status
\end{minipage} & \begin{minipage}[b]{\linewidth}\raggedright
Referência
\end{minipage} \\
\midrule\noalign{}
\endhead
\bottomrule\noalign{}
\endlastfoot
1 & Spazio \(\mathcal{H}\) degli stati egemonici & \Proved{} definição
operativa & Geometria \S{}2 \\
2 & Forma di \(Q_i\) (identica a \(Q_0\) Topografia) & \Proved{}
verificaçãodo & Geometria \S{}2 / DEQ$^2$ Cap.~3 \\
3 & Payoff \(\Pi_i\) identico a $\Pi$ Topografia & \Proved{} verificaçãodo
(Cap.~3) & Geometria \S{}3 / DEQ$^2$ Cap.~3 \\
4 & Teorema Folk \(\delta^* = 0{,}50\) & \Proved{} derivato & Geometria
\S{}3 \\
5 & \(T_s^{(1)}\) come soglia composita & \Hypoth{} hipótese operativa &
Memoriale DEQ \\
6 & EDSD come dinamica stocastica & \Proved{} formulazione standard
(Itô) & Geometria \S{}12-13 \\
7 & \(\varepsilon_{\text{dis}}\) ancorata a Bandura/Milgram & \Proved{}
ancoraggio empirico documentato & Memoriale DEQ \\
8 & Limite della DEQ$^1$ in assenza di \(A, \Omega\) & \Conj{} critica
costruttiva, motivante per DEQ$^2$ & DEQ$^2$ Cap.~3 \\
\end{longtable}}
}


